\documentclass[11pt]{article}
\usepackage[letterpaper,margin=1in]{geometry}
\usepackage[T1]{fontenc}
\usepackage{lmodern,microtype}
\usepackage{amsmath,amssymb,amsthm,mathtools}
\usepackage[hidelinks]{hyperref}
\numberwithin{equation}{section}
\newtheorem{theorem}{Theorem}[section]
\newtheorem{lemma}[theorem]{Lemma}
\newtheorem{proposition}[theorem]{Proposition}
\newtheorem{corollary}[theorem]{Corollary}
\theoremstyle{remark}\newtheorem{remark}[theorem]{Remark}
\DeclareMathOperator{\Ber}{Ber}
\DeclareMathOperator{\Bin}{Bin}
\DeclareMathOperator{\TV}{TV}
\newcommand{\PP}{\mathbb P}
\newcommand{\EE}{\mathbb E}
\newcommand{\one}{\mathbf 1}
\newcommand{\dd}{\,\mathrm d}
\newcommand{\Bcal}{\mathcal B}
\newcommand{\Dcal}{\mathcal D}
\title{The exact Poisson--binomial comparison\\in every dimension}
\author{Christopher D. Long\\
  \small\href{mailto:galizur@gmail.com}{\texttt{galizur@gmail.com}}}
\date{September 13, 2026}
\hypersetup{
  pdftitle={The exact Poisson-binomial comparison in every dimension},
  pdfauthor={Christopher D. Long}
}
\begin{document}
\maketitle
\begin{abstract}
For coordinatewise ordered Bernoulli parameters with a prescribed total
mean gap, the maximum difference of count tails is minimized by homogeneous
parameters. For a nondegenerate gap $0<\Delta<n$, in odd dimension
$n=2m+1\ge3$ the two minimizers are the homogeneous pairs whose likelihood
ratios equal one at count $m$ or $m+1$; in even dimension $n\ge4$ the
unique minimizer is the complementary homogeneous pair. We prove the comparison and these equality statements by a global
stationary-point reduction, with all boundary faces included. A normalized
switch-mass identity orders the homogeneous switching values. The same
extremizers give the exact product total-variation minimum. The exceptional
two-dimensional equality families are classified in the appendix.
\end{abstract}

\section{Statement}
For $p\in[0,1]^n$, let $S_p$ be a sum of independent Bernoulli variables
with parameters $p_i$. Put
\[
 D_k(p,q)=\PP(S_p\ge k)-\PP(S_q\ge k),\qquad
 T(p,q)=\max_{1\le k\le n}D_k(p,q),\qquad
 \Delta=\sum_i(p_i-q_i).
\]
For probability measures on a finite set, we use
$\TV(P,Q)=\tfrac12\sum_x|P(x)-Q(x)|$.
We assume $p\ge q$ coordinatewise. For $n\ge2$, $1\le j\le n-1$, and
$0<\gamma<1$, there is a unique pair $0<b<a<1$ satisfying
\begin{equation}\label{eq:switch}
 a-b=\gamma,\qquad a^j(1-a)^{n-j}=b^j(1-b)^{n-j}.
\end{equation}
Define
\[
 F_{n,j}(\gamma)=\PP(\Bin(n,a)\ge j)-\PP(\Bin(n,b)\ge j),
 \qquad
 \Bcal_n(\Delta)=F_{n,\lfloor n/2\rfloor}(\Delta/n).
\]
Set $\Bcal_n(0)=0$ and $\Bcal_n(n)=1$.

\begin{theorem}\label{thm:main}
For every $n\ge2$ and every $p,q\in[0,1]^n$ with $p\ge q$,
\begin{equation}\label{eq:main}
 T(p,q)\ge\Bcal_n(\Delta).
\end{equation}
For $n\ge3$ and $0<\Delta<n$, equality forces both vectors to be
homogeneous. More precisely:
\begin{enumerate}
\item If $n$ is even, equality holds exactly when
$p_i=(1+\Delta/n)/2$ and $q_i=(1-\Delta/n)/2$ for all $i$.
\item If $n=2m+1$ is odd, let $(a,b)$ solve \eqref{eq:switch} with
$j=m$ and $\gamma=\Delta/n$. Equality holds exactly for
\[
 (p,q)=(a\one,b\one)
 \quad\text{or}\quad
 (p,q)=((1-b)\one,(1-a)\one).
\]
\end{enumerate}
The exact minimum of
$\TV(\bigotimes_i\Ber(p_i),\bigotimes_i\Ber(q_i))$ under the same
constraint is also $\Bcal_n(\Delta)$, with the same equality cases for
$n\ge3$ and $0<\Delta<n$.
\end{theorem}
The exceptional equality cases for $n=2$ are classified in
Proposition~\ref{prop:n2-equality} of Appendix~\ref{app:n2}.

For odd $n=2m+1$,
the benchmark can also be written
\begin{equation}\label{eq:odd-integral}
 \Bcal_{2m+1}(\Delta)=(2m+1)\binom{2m}{m}
 \int_b^a [u(1-u)]^m\dd u,
\end{equation}
where $(a,b)$ is the switch with $j=m$. Equivalently, let $z>1$ be the
unique solution of
\[
 \frac{(z^m-1)(z^{m+1}-1)}{z^{2m+1}-1}=\frac{\Delta}{2m+1}.
\]
Then
\[
 b=\frac{z^m-1}{z^{2m+1}-1},\qquad
 a=\frac{z^{m+1}(z^m-1)}{z^{2m+1}-1}.
\]
These formulas follow directly from the equal-height equation.


The even-dimensional comparison and related product optimization were posed
by Kontorovich \cite{KontorovichMO,KontorovichTV}. The discussion of
\cite{KontorovichMO} includes Steinke's randomized symmetrization and
Sawin's fixed-mean tail-extremization approach. General lower bounds for
product total variation in terms of marginal distances are developed in
\cite{KontorovichTensorization}. More directly,
\cite[Theorem~1.5]{KontorovichHomogenization} compares Bernoulli products
with their mean-preserving homogenizations up to a universal constant.
Here only the total gap is fixed, and we determine the exact minimum and
all equality cases. The proof combines the deletion-mixture bound in
Lemma~\ref{lem:atom}, the switch identities in
Lemma~\ref{lem:switch-identities}, and a global extremizer argument.
The even-dimensional benchmark is recovered as
\[
 \Bcal_n(\Delta)=2^{-n}\binom n{n/2}
 \int_0^\Delta(1-u^2/n^2)^{n/2-1}\dd u.
\]
For $n=3$, \eqref{eq:odd-integral} becomes
\[
 \Bcal_3(\Delta)=\frac{2\Delta}{9}
 \left(1+\sqrt{1-\Delta^2/12}\right).
\]
The proof proceeds by induction from $n=2$.

\section{Bernoulli facts and the randomized test}\label{sec:bernoulli}
We use the likelihood-ratio order only for Bernoulli sums. Their mass
functions have interval support and are strictly log-concave at positive
interior triples. A deterministic shift preserves these facts. To see
strict log-concavity, remove the deterministic variables and write the
remaining probability generating polynomial as
$\prod_{i=1}^M(1-t_i+t_i z)$, where $0<t_i<1$. Its zeros are real and
negative. If its coefficients are $a_0,\ldots,a_M$, Newton's inequality
(see \cite{Pitman}) gives, for $1\le j\le M-1$,
\[
 a_j^2\ge a_{j-1}a_{j+1}
 \frac{(j+1)(M-j+1)}{j(M-j)}>a_{j-1}a_{j+1}.
\]

Here and below a comparison of ratios is understood by cross
multiplication when an endpoint mass is zero.
\begin{lemma}[Increasing a parameter]\label{lem:lr}
Increasing Bernoulli parameters, or adjoining a Bernoulli variable,
increases the law in likelihood-ratio order. If at least one parameter
increases strictly, the adjacent likelihood-ratio comparison is strict
when both adjacent masses are positive in both laws. In particular, if
$\PP(Z=k-1)=\PP(Z=k)>0$, and $Z'$ is obtained by a strict parameter
increase, then
\[
 \PP(Z'=k)>0\quad\Longrightarrow\quad
 \PP(Z'=k-1)<\PP(Z'=k).
\]
If the mass at $k$ disappears because the lower endpoint of the support
moves past $k$, both displayed masses of $Z'$ are zero instead.
On a parameter interval where both adjacent masses remain positive,
their difference can cross zero only from positive to negative.
\end{lemma}
\begin{proof}
With all other variables summed into $R$, changing a parameter from $u$
to $v>u$ changes the mass at $j$ from
$(1-u)r_j+ur_{j-1}$ to $(1-v)r_j+vr_{j-1}$. Their ratio increases with
$r_{j-1}/r_j$, which increases strictly in $j$ on the positive support by
strict log-concavity. Equivalently, the relevant adjacent determinant is
a positive multiple of $r_{j-1}^2-r_{j-2}r_j$.
Successive coordinate changes prove the assertion. Adjoining a variable
is the change of an appended parameter from zero to its new value.
\end{proof}

For ordered $p,q$ with $\Delta>0$ and $T(p,q)<1$, Lemma~\ref{lem:lr}
shows that the maximizing thresholds consist of one threshold or two
adjacent thresholds. Indeed the mass difference changes sign once, and
can vanish at at most one count in the overlapping support. Disjoint
supports would give $T=1$.

For $1\le k\le n-1$ and $0\le\lambda\le1$, define
\[
 H(x)=\lambda\PP(S_x\ge k)+(1-\lambda)\PP(S_x\ge k+1).
\]
An independent $E\sim\Ber(\lambda)$ gives
$H(x)=\PP(S_x+E\ge k+1)$. Write $H_i=\partial_i H$ and
$C_{ij}=\partial_i\partial_j H$. Direct differentiation gives
\begin{align}
 H_i(x)&=\PP(S_{x,-i}+E=k),\label{eq:gradient}\\
 C_{ij}(x)&=\PP(S_{x,-ij}+E=k-1)-\PP(S_{x,-ij}+E=k),\label{eq:coefficient}\\
 H_i(x)-H_j(x)&=(x_j-x_i)C_{ij}(x).\label{eq:grad-difference}
\end{align}
A pure threshold is included by taking $\lambda=0$ or $1$.
For fixed other coordinates, splitting two equal coordinates $u,u$ into
$u+t,u-t$ changes $H$ by exactly $-C_{ij}t^2$.

\begin{lemma}[A maximal-atom bound]\label{lem:atom}
Put
\[
 \kappa_n=\binom n{\lfloor n/2\rfloor}
 \left(\frac{\lfloor n/2\rfloor}{n}\right)^{\lfloor n/2\rfloor}
 \left(\frac{\lceil n/2\rceil}{n}\right)^{\lceil n/2\rceil}.
\]
Every $(n-1)$-trial Bernoulli sum has a mass at least $\kappa_n$.
If at most two of the gaps $p_i-q_i$ are positive, then
\begin{equation}\label{eq:two-gaps}
 T(p,q)\ge\kappa_n\Delta.
\end{equation}
\end{lemma}
\begin{proof}
First, an $n$-trial Bernoulli sum of integer mean has probability at
that integer at least $\kappa_n$. To prove this, minimize that probability
at fixed mean. The fixed-mean reduction to equal interior parameters is
classical; see Hoeffding \cite{Hoeffding}. Here is the two-coordinate
argument needed in this case. Among the minimizers choose one with least
sum of squared parameters. If two interior parameters $u\ne v$ occur,
then, with the others fixed, the objective has the form
$A+B(u+v)+Cuv$. The transfer $(u,v)\mapsto(u+t,v-t)$ is feasible for
both signs of sufficiently small $t$, so minimality gives
$C(v-u)=0$. Thus $C=0$, and replacing $u,v$ by their average preserves
the objective while strictly decreasing the sum of squares, a
contradiction. Consequently an extremizer is a deterministic shift plus
$\Bin(M,\ell/M)$ with integer $\ell$ and $M\le n$.
For
\[
 b_{M,\ell}=\binom M\ell(\ell/M)^\ell(1-\ell/M)^{M-\ell},
 \qquad h(x)=x\log(1+1/x),\quad h(0)=0,
\]
ratio calculations give
\[
 \log\frac{b_{M,\ell}}{b_{M,\ell+1}}
 =h(M-\ell-1)-h(\ell),\qquad
 \log\frac{b_{M+1,\ell}}{b_{M,\ell}}=h(M-\ell)-h(M).
\]
The increasing function $h$ shows that the least $n$-trial integer-mean
mass is $\kappa_n$, and that increasing $M$ with $\ell$ fixed can only
decrease the mass. Deterministic cases give mass one.

If $\EE Z=k+f$ with $0\le f\le1$, adjoin
$\Ber(1-f)$. Its mass at the integer mean $k+1$ is
$(1-f)\PP(Z=k)+f\PP(Z=k+1)\ge\kappa_n$, proving the first assertion.
This extra Bernoulli variable explains the index $n$ in the bound for an
$(n-1)$-trial sum.

For the second assertion, the case $\Delta=0$ is immediate; assume
$\Delta>0$. Put $r_i=(p_i+q_i)/2$, $d_i=p_i-q_i$ and
$w_i=d_i/\Delta$. When at most two $d_i$ are nonzero, symmetric expansion
of the generating polynomials gives exactly
\[
 \frac{D_{j+1}}\Delta
 =\sum_iw_i\PP(S_{r,-i}=j).
\]
This deletion mixture is an $(n-1)$-trial Bernoulli law. For completeness,
its generating polynomial is
$G(z)=\sum_iw_i\prod_{\ell\ne i}(1-r_\ell+r_\ell z)$.
For distinct interior $r_i$ and positive weights, factor out the full
product: the remaining rational function has positive residues at its
negative poles and therefore one zero between each successive pair of
poles. Thus $G$ has only negative real zeros, and $G(1)=1$ factors it
into Bernoulli generating factors. Continuity covers all degeneracies.
The first assertion now proves \eqref{eq:two-gaps}.
\end{proof}

We will repeatedly use one more elementary consequence of
Lemma~\ref{lem:lr}. If $W$ is a fixed Bernoulli sum and
\begin{equation}\label{eq:gM}
 g_M(t)=\PP(\Bin(M,t)+W=k),
\end{equation}
then
\[
 g_M'(t)=M\{\PP(\Bin(M-1,t)+W=k-1)
                 -\PP(\Bin(M-1,t)+W=k)\}.
\]
The derivative formula is used for $M\ge1$; $g_0$ is simply a
constant. On $0<t<1$ the support is fixed. Thus the derivative changes
sign at most once, from positive to negative. An interior critical point
of positive height is its maximum; it is strict when $M\ge2$. For
$M=1$ an interior critical point means that $g_1$ is constant.

We also need the following adjoining rule, with its positivity condition
explicit: if a Bernoulli sum $Z$ satisfies
\[
 \PP(Z=k-1)\le\PP(Z=k),\qquad \PP(Z=k)>0,
\]
then adjoining further independent Bernoulli variables cannot increase
its mass at $k$. Indeed, write $a=\PP(Z=k-2)$, $b=\PP(Z=k-1)$ and
$c=\PP(Z=k)>0$. Log-concavity and interval support give
$a\le b\le c$. After adjoining $\Ber(t)$ the relevant masses are
$(1-t)b+ta$ and $(1-t)c+tb$, so their order persists and the mass at
$k$ does not increase. If that mass becomes zero, the support has moved
strictly above $k$, and further adjoining cannot recreate it.
If $b=c>0$ and $t>0$, strict log-concavity gives $a<b$; the first
addition makes the adjacent difference strictly negative, and a second
addition with positive parameter strictly decreases the mass at $k$.
All uses of this rule below start with a positive mass at $k$.

\begin{corollary}[Derivative criterion for adjoining]\label{cor:adjoining}
Let $g_M$ be as in \eqref{eq:gM}, with $M\ge1$ and $0<t<1$.
If $g_M(t)>0$ and $g_M'(t)\le0$, then
\[
 \PP(\Bin(M,t)+W=k-1)\le\PP(\Bin(M,t)+W=k).
\]
Consequently adjoining further independent Bernoulli variables cannot
increase the mass at $k$.
\end{corollary}
\begin{proof}
Write $Y=\Bin(M-1,t)+W$ and $y_j=\PP(Y=j)$. The derivative condition
gives $y_{k-1}\le y_k$, and $g_M(t)>0$ then implies $y_k>0$.
Log-concavity and interval support give $y_{k-2}\le y_{k-1}\le y_k$.
For $Z=Y+\Ber(t)$, it follows that
\[
 \PP(Z=k-1)=(1-t)y_{k-1}+ty_{k-2}
 \le(1-t)y_k+ty_{k-1}=\PP(Z=k).
\]
The adjoining rule applies.
\end{proof}

\section{The homogeneous problem in both parities}
Suppress the first index in $F_{n,j}$, and write $c=j/n$.
Let $(a,b)$ satisfy \eqref{eq:switch}, and set
\[
 v=c(1-c),\quad x=a-c,\quad y=c-b,\quad
 R=\frac{xy}{v},\quad f_j(\gamma)=\binom nj a^j(1-a)^{n-j}.
\]
Here $x,y>0$, and
\begin{equation}\label{eq:Rden}
 v(1-R)=c(1-a)(1-b)+(1-c)ab>0.
\end{equation}
The binomial masses at count $j$ agree, so thresholds $j,j+1$ tie and
both equal $F_j(\gamma)$.

\begin{lemma}[Switch identities]\label{lem:switch-identities}
Primes denote differentiation in $\gamma$. Then
\begin{equation}\label{eq:switch-derivatives}
 F_j'=\frac{nf_j}{1-R},\qquad
 \frac{f_j'}{f_j}=-\frac{nR}{\gamma(1-R)}.
\end{equation}
Consequently,
\begin{equation}\label{eq:switch-normalization}
 F_j(\gamma)=(n+1)\int_0^\gamma f_j(u)\dd u-\gamma f_j(\gamma),\qquad
 \int_0^1f_j(u)\dd u=\frac1{n+1}.
\end{equation}
For fixed $\gamma\in(0,1)$, $R$ is strictly decreasing in $c$ on
$(0,1/2)$. Also $R'\le2R/\gamma$.
\end{lemma}
\begin{proof}
The equation defining $b$ has derivative
\[
 -\gamma\left\{\frac c{ab}+\frac{1-c}{(1-a)(1-b)}\right\}<0
\]
in $b$, with opposite infinite endpoint limits. This proves existence
and uniqueness of the switch. Differentiating the equal-height equation
and $a-b=\gamma$, and using \eqref{eq:Rden}, gives
\[
 a'=\frac{y a(1-a)}{\gamma v(1-R)},\qquad
 b'=-\frac{x b(1-b)}{\gamma v(1-R)}.
\]
Differentiation of the binomial tail proves
\eqref{eq:switch-derivatives}. In particular $F_j'=nf_j-\gamma f_j'$.
Integration, followed by $F_j(1)=1$ and $f_j(1)=0$, proves
\eqref{eq:switch-normalization}.

We give the geometric monotonicity explicitly. Suppose $c\le1/2$, and
put $z=1-2c$ and $e=a+b-2c=x-y$. Let $A=\operatorname{logit}(a)$,
$B=\operatorname{logit}(b)$, $M=(A+B)/2$, $s=(A-B)/2$, and
$\sigma(t)=e^t/(1+e^t)$. The equal-height equation says
\[
 c=\frac1{2s}\int_{M-s}^{M+s}\sigma(t)\dd t.
\]
Thus $M\le0$. The identity
\[
 \sigma(M+u)+\sigma(M-u)
 =1+\frac{\sinh M}{\cosh M+\cosh u}
\]
shows that the paired average increases with $u\ge0$. Hence
$\sigma(M)\le c\le(a+b)/2$, strictly if $c<1/2$. The first inequality
is equivalent to
\[
 c^2(1-a)(1-b)-(1-c)^2ab=v(zR-e)\ge0,
\]
and the second gives $e\ge0$.

Now keep $\gamma$ fixed and differentiate in the continuous parameter
$c$. The derivative $t=\partial_c a=\partial_c b$ is positive, since
the derivative of the switch equation in $c$ is $A-B>0$ whereas its
derivative in $b$ is negative. Therefore
\begin{equation}\label{eq:Rcentering}
 \partial_cR=-\frac{(zR-e)+et}{v}<0\qquad(0<c<1/2).
\end{equation}
Finally set $\theta=x/\gamma$. Differentiation gives
$\theta'=(zR-e)/[\gamma^2(1-R)]\ge0$, and $e\ge0$ gives
$\theta\ge1/2$. Since
$R/\gamma^2=\theta(1-\theta)/v$, its derivative is nonpositive.
Reflection covers $c>1/2$ and proves the last assertion.
\end{proof}

\begin{proposition}[Strict ordering of switches]\label{prop:switch-order}
For $1\le i<j\le n/2$ and $0<\gamma<1$,
\[
 F_{n,i}(\gamma)>F_{n,j}(\gamma).
\]
Consequently the homogeneous minimum is $\Bcal_n(\Delta)$, with exactly
the central switch when $n$ is even and the two reflected central
switches when $n$ is odd.
For $n\ge3$, $\Bcal_n$ is strictly concave on $(0,n)$ and
\begin{equation}\label{eq:strict-linear}
 0<\Bcal_n(\Delta)<\min\{\kappa_n\Delta,1\}
 \qquad(0<\Delta<n).
\end{equation}
\end{proposition}
\begin{proof}
By \eqref{eq:Rcentering}, $R_i>R_j$. Thus
\[
 \left(\log\frac{f_i}{f_j}\right)'
 =-\frac n\gamma\left(\frac{R_i}{1-R_i}-\frac{R_j}{1-R_j}\right)<0.
\]
The equal integrals in \eqref{eq:switch-normalization} imply that
$h=f_i-f_j$ changes sign exactly once, from positive to negative.
Where $h>0$, writing $A_\ell=nR_\ell/[\gamma(1-R_\ell)]$ gives
$h'=-A_i h-(A_i-A_j)f_j<0$. Before the crossing,
$\int_0^\gamma h\ge\gamma h(\gamma)>0$, and afterwards
$\int_0^\gamma h=-\int_\gamma^1h>0$. In both cases
\[
 F_i-F_j=(n+1)\int_0^\gamma h-\gamma h(\gamma)>0.
\]
Reflection gives $F_{n,j}=F_{n,n-j}$.

It remains to justify that switches exhaust the homogeneous minima.
Define
\[
 \Dcal_h(a,b)=D_h(a\one,b\one),\qquad
 M_\gamma(b)=\max_{1\le h\le n}\Dcal_h(b+\gamma,b).
\]
At fixed $\gamma$, each threshold difference has the form
\[
 \Dcal_h(b+\gamma,b)=n\int_b^{b+\gamma}
 \binom{n-1}{h-1}u^{h-1}(1-u)^{n-h}\dd u.
\]
Its derivative has the sign of the ratio of the two endpoint densities
minus one. Strict log-concavity makes that ratio strictly decreasing in
$b$, so the individual threshold difference has no interior local minimum.
Suppose that $b_0\in(0,1-\gamma)$ is a local minimum of $M_\gamma$.
If exactly one threshold is active at $b_0$, its strict separation from
the other finitely many thresholds persists in a neighborhood. There
$M_\gamma$ equals that threshold difference, a contradiction. Thus at
least two thresholds are active. Here $0<b_0<b_0+\gamma<1$, so both
count laws have full support and $T<1$. The active-threshold property in
Section~\ref{sec:bernoulli} therefore shows that they are exactly two
adjacent thresholds
$k,k+1$. Their equality is equivalent to
\[
 \PP(\Bin(n,b_0+\gamma)=k)=\PP(\Bin(n,b_0)=k),
\]
so $b_0$ is a switch. The homogeneous minimum is therefore attained at
a switch or at $b=0,1-\gamma$.
At an endpoint its value is $1-(1-\gamma)^n$, strictly larger than a
central switch value: in the coordinatewise uniform coupling of a
central homogeneous pair, the event of at least one discrepancy has
this probability, and the positive-probability event of exactly one
discrepancy with all other coordinates zero does not cross threshold
$\lfloor n/2\rfloor+1$.

Finally \eqref{eq:switch-derivatives} and $R'\le2R/\gamma$ give
\[
 (\log F_j')'\le-\frac{(n-2)R}{\gamma(1-R)}<0\quad(n\ge3).
\]
For the central index $j=\lfloor n/2\rfloor$, as $\gamma\downarrow0$
we have $a,b\to j/n$ and $R\to0$. Since $\Delta=n\gamma$,
\[
 \Bcal_n'(0+)
 =\lim_{\gamma\downarrow0}\frac1n F_{n,j}'(\gamma)
 =\lim_{\gamma\downarrow0}\frac{f_j(\gamma)}{1-R(\gamma)}
 =\binom nj\left(\frac jn\right)^j
            \left(1-\frac jn\right)^{n-j}
 =\kappa_n.
\]
Strict concavity, together with $\Bcal_n(0)=0$ and $\Bcal_n(n)=1$,
proves \eqref{eq:strict-linear}.
\end{proof}

\section{Removing a common deterministic coordinate}
The following strict dimension improvement supplies the boundary step
in the induction.
\begin{lemma}\label{lem:dimension}
Let $N\ge2$. For every $0<\Delta\le N$, there is an ordered
$(N+1)$-dimensional pair of gap
$\Delta$ with tail objective strictly less than $\Bcal_N(\Delta)$.
\end{lemma}
\begin{proof}
For $\Delta=N$, any strictly interior ordered pair of that gap has
$T<1=\Bcal_N(N)$. Suppose $0<\Delta<N$ and take a homogeneous central
switch pair $(a\one,b\one)$ in dimension $N$, with switch index $j$,
common mass $f>0$, and geometric quantity $R>0$ from
Lemma~\ref{lem:switch-identities}. Set
\[
 \rho=\frac{(N-j)ab}{(N-j)ab+j(1-a)(1-b)}\in(0,1).
\]
The density
\[
 g(u)=\rho\PP(\Bin(N-1,u)=j-1)
       +(1-\rho)\PP(\Bin(N-1,u)=j)
\]
satisfies, by substitution,
\begin{equation}\label{eq:rho-density}
 g(a)=g(b)=\frac f{1-R}.
\end{equation}
Adjoin the common parameter $\rho$. To identify the active threshold,
write $D_h$ for the old tail differences and set $D_0=D_{N+1}=0$.
The new differences satisfy
\[
 \widetilde D_h=(1-\rho)D_h+\rho D_{h-1},\qquad 1\le h\le N+1.
\]
Exactly the old thresholds $j,j+1$ attain $\Bcal_N(\Delta)$.
Since $0<\rho<1$, the new threshold $j+1$ still has this value,
whereas every other new threshold has a strictly smaller value: its
convex combination includes at least one old value below the maximum.
Thus $j+1$ is the unique maximizing threshold.
For small $\varepsilon>0$ replace the parameters by
\[
 p_\varepsilon=
 \left((a-\varepsilon/(2N))^N,\rho+\varepsilon/2\right),\qquad
 q_\varepsilon=
 \left((b+\varepsilon/(2N))^N,\rho-\varepsilon/2\right).
\]
Here a superscript on a coordinate denotes repetition. The gap remains
$\Delta$ and the pair is feasible. By \eqref{eq:rho-density}, the unique
active tail has derivative at zero
\[
 f-\frac{g(a)+g(b)}2=f-\frac f{1-R}<0.
\]
It remains the maximizing tail for sufficiently small $\varepsilon$.
This proves the strict improvement.
\end{proof}

\section{The global minimizer}
We prove Theorem~\ref{thm:main} by induction. For $n=2$,
$D_1+D_2=\Delta$, so $T\ge\Delta/2=\Bcal_2(\Delta)$, attained by the
complementary homogeneous pair. Fix $n\ge3$, assume the theorem's
minimum-value assertion in dimension $n-1$, and let $(p,q)$ be a global
minimizer at a fixed gap $0<\Delta<n$. Compactness gives its existence.
Its value is at most the homogeneous candidate $\Bcal_n(\Delta)$.

By Lemma~\ref{lem:atom} and \eqref{eq:strict-linear}, at least three
gaps are positive. By Lemma~\ref{lem:dimension}, no coordinate can be
common deterministic. Such a coordinate requires $\Delta\le n-1$;
deleting it would give
$T\ge\Bcal_{n-1}(\Delta)$, whereas a strict improvement is available.
Thus $p_i>0$ and $q_i<1$ for every $i$. Also $p$ cannot be identically
one, nor $q$ identically zero. For example, if $p=\one$, then
\[
 T=1-\prod_iq_i\ge1-(1-\Delta/n)^n>\Bcal_n(\Delta),
\]
where the strict endpoint comparison was proved above.

Write
\[
 I=\{i:p_i>q_i\},\qquad C=\{i:p_i=q_i\},\qquad
 r=|I|\ge3,\quad s=|C|.
\]
All common parameters, when present, are strictly interior.

\subsection{Stationarity and common coordinates}
There are one or two adjacent active thresholds. Choose their convex
combination $H$ as in \eqref{eq:gradient} so that the first-order
optimality conditions hold. Existence follows from the separating
hyperplane theorem applied to the feasible tangent cone of the linear
constraints and the active gradients. With a multiplier $c$ for the
total gap, these conditions read, for $i\in I$,
\begin{align}
 p_i<1&\Longrightarrow H_i(p)=c,&p_i=1&\Longrightarrow H_i(p)\le c,
 \label{eq:KKTp}\\
 q_i>0&\Longrightarrow H_i(q)=c,&q_i=0&\Longrightarrow H_i(q)\le c,
 \label{eq:KKTq}
\end{align}
and, for $i\in C$,
\begin{equation}\label{eq:KKTC}
 H_i(p)=H_i(q)=c+\nu_i,\qquad \nu_i\ge0.
\end{equation}
The signs follow, for example, by using
$H(p)-H(q)-c(\sum_i(p_i-q_i)-\Delta)-\sum_{i\in C}\nu_i(p_i-q_i)$
as the Lagrangian, together with the endpoint constraints.

We first establish, without any nonconstancy assumption, that
\begin{equation}\label{eq:c-positive}
 c>0.
\end{equation}
Write the supports of $S_p,S_q$ as $[L,n]$ and $[0,U]$.
Since $T<1$, we have $L\le U$. For every changing coordinate $i$,
the derivative of the pure threshold $h$ at $p$ is positive whenever
$L+1\le h\le n$, and its derivative at $q$ is positive whenever
$1\le h\le U$. This follows from the deletion supports and includes
$p_i=1$ and $q_i=0$. Every threshold belongs to at least one of these
two ranges. Therefore their active convex combination satisfies
\[
 0<H_i(p)+H_i(q)\le2c\qquad(i\in I),
\]
where the upper bound follows from
\eqref{eq:KKTp}--\eqref{eq:KKTq}. This proves \eqref{eq:c-positive}.

If there is only one active threshold $h$ and a common coordinate
$t$, stationarity in $t$ gives
$\PP(S_{p,-i}=h-1)=\PP(S_{q,-i}=h-1)$.
These are positive masses. Indeed, write the original supports as
$[L,n]$ and $[0,U]$. Deleting the common interior coordinate gives
supports $[L,n-1]$ and $[0,U-1]$. Both masses at $h-1$ could vanish
only if $h\le L$ and $h\ge U+1$, forcing $L>U$ and hence $T=1$.
Likelihood-ratio ordering then says that the two deletion tails at
$h-1,h$ both maximize their difference. Adjoining any common Bernoulli
parameter preserves this maximum. Taking that parameter to zero creates
a common deterministic coordinate at the same minimum, a contradiction.

Consequently, if $C\ne\varnothing$, there are two active thresholds
$k,k+1$. For a common coordinate $t$, put
$u_j=\PP(S_{p,-i}=j)-\PP(S_{q,-i}=j)$. The active tie and
\eqref{eq:KKTC} give
\[
 tu_{k-1}+(1-t)u_k=0,\qquad
 \lambda u_{k-1}+(1-\lambda)u_k=0.
\]
The tied count lies in the overlap, so $L\le k\le U$.
The deletion supports are $[L,n-1]$ and $[0,U-1]$; the upper deletion
mass at $k$ and the lower deletion mass at $k-1$ are positive. If
$u_{k-1}=u_k=0$, equality would make both adjacent masses positive
in both deletion laws. At least one changing coordinate remains, so
strict likelihood-ratio ordering would contradict the two equalities.
Thus the two adjacent differences cannot both vanish, and hence
\begin{equation}\label{eq:commonlambda}
 t=\lambda\in(0,1)\quad\text{for every common coordinate}.
\end{equation}
In particular all common coordinates agree. Their derivative is exactly
the original mass at the tied count. That mass is positive because
$S_p$ has support $[L,n]$ and $L\le k\le n-1$. Thus
\begin{equation}\label{eq:commonf}
 H_i(p)=H_i(q)=f:=\PP(S_p=k)=\PP(S_q=k)>0
 \quad(i\in C),\qquad f=c+\nu,\quad\nu\ge0.
\end{equation}

\subsection{One changing vector is homogeneous}
For $i,j\in I$, \eqref{eq:KKTp} and \eqref{eq:grad-difference} show
\[
 p_i\ne p_j\quad\Longrightarrow\quad C_{ij}(p)\ge0;
\]
indeed distinct interior entries have equal gradients, and a larger
entry equal to one has no larger gradient. Likewise
\[
 q_i\ne q_j\quad\Longrightarrow\quad C_{ij}(q)\le0.
\]
Suppose that both $p_i\ne p_j$ and $q_i\ne q_j$. The laws in
\eqref{eq:coefficient} are ordered in likelihood ratio, and at least
one of the other changing coordinates increases strictly because
$r\ge3$. The upper deletion law has positive mass at $k-1$: this follows
from a free gradient equal to $c>0$ and its nonnegative adjacent
difference. Similarly the lower deletion law has positive mass at $k$.
Their ordered interval supports therefore make both adjacent masses
positive in both laws. If $A_p,B_p$ denote the upper deletion masses
at $k-1,k$, and $A_q,B_q$ the corresponding lower deletion masses,
then the coefficient signs and Lemma~\ref{lem:lr} would give
\[
 1\le\frac{A_p}{B_p}<\frac{A_q}{B_q}\le1,
\]
a contradiction. Thus every pair satisfies $p_i=p_j$ or $q_i=q_j$.
The elementary consequence is that one of $p_I,q_I$ is constant: two
different values in one vector force all entries of the other to agree.
Reflection lets us assume
\begin{equation}\label{eq:pa}
 p_i=a\quad(i\in I).
\end{equation}
Here reflection means $(p,q)\mapsto(\one-q,\one-p)$; it preserves
$\Delta$ and $T$, with $D_h$ sent to $D_{n-h+1}$.

If $a=1$, then $s>0$, since the all-one vector was excluded. By
\eqref{eq:commonf} and \eqref{eq:KKTp}, for a common coordinate $i$
and a changing coordinate $j$,
\[
 (1-\lambda)C_{ij}(p)=H_i(p)-H_j(p)\ge0.
\]
Consider
$h(t_1,\ldots,t_r)=\PP(\sum_{j=1}^r\Ber(t_j)+\Bin(s,\lambda)=k)$.
Its partial derivatives at $t=\one$ equal $C_{ij}(p)\ge0$.
Along the path from $q_I$ to $\one$, the deletion laws are strictly
smaller in likelihood ratio than the endpoint deletion law
$(r-1)+\Bin(s,\lambda)$. Its two relevant masses are positive because
$r\le k\le n-1$, so $k-r$ and $k-r+1$ both lie in $\{0,\ldots,s\}$. All partial derivatives are therefore positive in the
interior of this path. It follows that
$h(\one)>h(q_I)$, contradicting the active mass equality. Hence
\begin{equation}\label{eq:a-interior}
 0<a<1.
\end{equation}

\subsection{The other changing vector is homogeneous}
Fix the common coordinates and set
\[
 W\sim\Bin(s+1,\lambda),\qquad
 g_M(t)=\PP(\Bin(M,t)+W=k).
\]
This notation includes deterministic $W$ when $s=0$ and $\lambda$ is
an endpoint. By \eqref{eq:pa}, \eqref{eq:a-interior}, and
\eqref{eq:KKTp},
\begin{equation}\label{eq:ca}
 c=g_{r-1}(a)>0.
\end{equation}
Every positive $q_i$, $i\in I$, is less than $a$ and is a free coordinate.
Its gradient is $c$.

There are at most two distinct positive values among these $q_i$.
Indeed, suppose three such coordinates have distinct values
$t_1,t_2,t_3$, and let $r_j$ be the masses of the law with those three
coordinates deleted and $E$ adjoined. Put
$A=r_{k-2}-r_{k-1}$ and $B=r_{k-1}-r_k$.
The pair coefficients vanish by \eqref{eq:grad-difference}, giving
$(1-t_\ell)B+t_\ell A=0$ for $\ell=1,2,3$.
Subtracting two equations yields $A=B$, and then $A=B=0$.
Thus $r_{k-2}=r_{k-1}=r_k$. Their common value equals a free gradient,
which is a convex combination of these three masses, and hence equals
$c>0$. This contradicts strict log-concavity.

If two positive values $u<v$ occur, the smaller value occurs only once.
Otherwise split two $u$ coordinates. The coefficient for a pair $u,v$
is zero, so its deletion law $Z_{uv}$, with $E$ adjoined, has equal
masses at $k-1,k$. The free-gradient identity
\[
 c=(1-v)\PP(Z_{uv}=k)+v\PP(Z_{uv}=k-1)
\]
shows that both equal masses are $c>0$. The deletion law for two $u$
coordinates replaces the remaining $u$ by $v$. Since $0<u<v<a<1$,
this replacement preserves the support and both adjacent masses stay
positive. Strict likelihood-ratio ordering therefore gives $C_{uu}<0$. The objective $H(p)-H(q)$ then
has negative second variation $C_{uu}t^2$ along the split.
For a single active threshold this is an immediate contradiction.
For two active thresholds keep their tie exactly: an unequal $u,v$
transfer has nonzero first derivative of
$G=\PP(S_p=k)-\PP(S_q=k)$, whereas the equal-coordinate split has zero
first derivative. To check nonvanishing, write the coefficient as
$\lambda A+(1-\lambda)B=0$, where
\[
 A=\PP(R=k-2)-\PP(R=k-1),\qquad
 B=\PP(R=k-1)-\PP(R=k)
\]
and $R$ excludes those two coordinates, without $E$. If $A-B=0$,
then $A=B=0$, so the three consecutive masses of $R$ are equal.
Their common value is $H_i(q)=c>0$ for the coordinate with value $u$:
adjoining $\Ber(v)$ and $E$ takes a convex combination of those three
masses. Strict log-concavity gives a contradiction. In fact, along the
transfer that increases a $u$ coordinate and decreases a $v$ coordinate,
\[
 \frac{\dd G}{\dd s}=-(v-u)(A-B)\ne0.
\]
Let $v_0$ denote the equal-coordinate split direction and $w_0$ this
unequal transfer direction, both acting on $q$ alone. For
$q(t,s)=q+t v_0+s w_0$, we have $G_t(0,0)=0$ and $G_s(0,0)\ne0$.
The implicit-function theorem gives $s(t)=O(t^2)$ with $G(t,s(t))=0$.
Both directions have coordinate sum zero. Moreover, for
$J(t,s)=H(p)-H(q(t,s))$, the free-gradient equalities give
$J_t(0,0)=J_s(0,0)=0$. Consequently
\[
 J(t,s(t))=J(0,0)+C_{uu}t^2+O(t^3)<J(0,0)
\]
for sufficiently small nonzero $t$. All other thresholds remain
strictly below the two tied active thresholds by continuity. On this
curve $T=J$, so this is a genuine feasible decrease of the original
maximum, not merely of a supporting test. All moves stay strictly
between zero and $a$, contradicting minimality.

Suppose now that the positive entries are one $u<v$ and $M$ copies of
$v$, with $b$ zero entries; thus $r=M+b+1$. Deleting $u$ gives
$c=g_M(v)$, while the zero $u,v$ coefficient gives $g_M'(v)=0$.
Consequently $g_M(v)$ is its maximum. If $M\ge2$, then $a>v$
gives $g_M(a)<g_M(v)$ and $g_M'(a)<0$. The mass $g_M(a)$ is
positive because the support is fixed on $0<t<1$ and $g_M(v)=c>0$.
Corollary~\ref{cor:adjoining} therefore gives
\[
 g_{M+b}(a)\le g_M(a)<g_M(v)=c,
\]
contradicting \eqref{eq:ca}.
If $M=1$, the function is constant and $b\ge1$ because $r\ge3$.
The law $W$ has equal positive masses at $k-1,k$. Adding one
$\Ber(a)$ preserves the mass $c$ at $k$ but makes the adjacent
difference strictly negative. The second addition therefore gives
$g_2(a)<g_1(a)=c$, and further additions cannot increase the mass.
Thus there cannot be two positive values.

Finally suppose that there are $M$ copies of a single positive value
$v$ and $b\ge1$ zeros, so $r=M+b$. A positive-coordinate derivative
and a zero-coordinate derivative give
\[
 c=g_{M-1}(v),\qquad g_M(v)\le c.
\]
For $R_0=\Bin(M-1,v)+W$, the identity
\[
 g_M(v)-g_{M-1}(v)
 =v\{\PP(R_0=k-1)-\PP(R_0=k)\}\le0
\]
shows that $\PP(R_0=k-1)\le\PP(R_0=k)=c>0$, since $v>0$.
Adjoin $b$ Bernoulli variables with
parameter $a$ and then raise the original $M-1$ parameters from $v$
to $a$. Every step can only decrease the mass at $k$. For the raising
steps, its deletion law dominates $R_0$: one appended $a>v$ replaces
the deleted $v$, and any remaining appended variables are extra.
The decrease is strict. If the initial adjacent masses are unequal,
the first addition is strict. If they tie, two additions are strict
in aggregate; when $b=1$, there is a parameter to raise because
$r\ge3$, and that raising step is strict. Hence again
$g_{r-1}(a)<c$, a contradiction.

We have proved that $q_I=b\one$ for some $0\le b<a$. The case $b=0$
is excluded either by the all-zero-vector endpoint comparison when
$s=0$, or by reflection of the $a=1$ argument when $s>0$. Therefore
\begin{equation}\label{eq:ab-interior}
 0<b<a<1.
\end{equation}

\subsection{No unchanged coordinate remains}
Suppose $s>0$. Recalling \eqref{eq:commonlambda} and
\eqref{eq:ab-interior}, we have
\[
 p=(a^r,\lambda^s),\qquad q=(b^r,\lambda^s),\qquad 0<b<a<1.
\]
Set
\begin{align*}
 g(t)&=\PP(\Bin(r-1,t)+\Bin(s+1,\lambda)=k),\\
 h(t)&=\PP(\Bin(r-1,t)+\Bin(s,\lambda)=k-1)
       -\PP(\Bin(r-1,t)+\Bin(s,\lambda)=k).
\end{align*}
The free gradients and common-coordinate gradients give
\begin{equation}\label{eq:nu}
 g(a)=g(b)=c,\qquad
 \nu=(a-\lambda)h(a)=(b-\lambda)h(b)\ge0.
\end{equation}
Since $r-1\ge2$, $g$ has a unique mode $u\in(b,a)$.
At that point the law
$\Bin(r-2,u)+\Bin(s+1,\lambda)$ has equal positive masses at $k-1,k$.
All laws compared below have $n-1$ strictly interior Bernoulli
parameters, including those obtained by intermediate replacements.
They therefore have positive masses at both $k-1$ and $k$.

If $\lambda\le b$, then $u>\lambda$. Replacing one $\lambda$
parameter by $u$ increases this deletion law strictly in likelihood
ratio, so $h(u)<0$ and hence $h(a)<0$. This makes
$(a-\lambda)h(a)<0$. If $\lambda\ge a$, the reflected argument gives
$h(b)>0$ and $(b-\lambda)h(b)<0$. In the remaining case
$b<\lambda<a$, nonnegativity in \eqref{eq:nu} would force
$h(b)\le0\le h(a)$, impossible because increasing $t$ crosses the
adjacent-mass equality only from positive to negative, strictly.
Every case contradicts \eqref{eq:nu}. Thus $s=0$.

The global minimizer is therefore fully homogeneous, with strictly
interior parameters. Proposition~\ref{prop:switch-order} identifies it
as precisely the central switch or the two reflected central switches.
This proves \eqref{eq:main} and all equality cases for $n\ge3$, and
completes the induction.

\section{Product total variation and endpoints}
For the threshold event $E_k=\{x:\sum_i x_i\ge k\}$,
\[
 \TV\left(\bigotimes_i\Ber(p_i),\bigotimes_i\Ber(q_i)\right)
 \ge D_k(p,q).
\]
At a homogeneous pair the product likelihood ratio depends only on
the count and increases strictly with it, so product total variation
equals the maximizing tail difference. The sharp minimum is therefore
$\Bcal_n(\Delta)$, and equality for $n\ge3$ forces equality in the tail
comparison, giving the same classification.
At $\Delta=0$, ordering forces $p=q$, and every such pair minimizes.
At $\Delta=n$, feasibility forces $p=\one,q=0$.
The nonunique two-dimensional equality families are proved in
Appendix~\ref{app:n2}.

\appendix
\section{Two-dimensional equality families}\label{app:n2}
For $x\in[0,1]^2$, write $P_x=\Ber(x_1)\otimes\Ber(x_2)$.
\begin{proposition}\label{prop:n2-equality}
Let $p,q\in[0,1]^2$ satisfy $p\ge q$, and put
$d_i=p_i-q_i$, $\Delta=d_1+d_2$, and $A=p_1p_2-q_1q_2$. Then
\begin{align}
 T(p,q)&=\frac\Delta2+\left|A-\frac\Delta2\right|,
 \label{eq:n2-tail}\\
 \TV(P_p,P_q)&=\frac\Delta2+\frac12\bigl(|d_1-A|+|d_2-A|\bigr).
 \label{eq:n2-product}
\end{align}
For $0<\Delta<2$, both minima are $\Delta/2$. Tail equality holds
exactly when $A=\Delta/2$. Product equality holds exactly for
\begin{equation}\label{eq:n2-family}
 p=(r+d,1-r),\qquad q=(r,1-d-r),\qquad
 d=\Delta/2,\quad 0\le r\le1-d.
\end{equation}
The product-minimizing family is a proper subset of the tail-minimizing
family for every $0<\Delta<2$.
\end{proposition}
\begin{proof}
Ordering gives $0\le A=d_1p_2+d_2q_1\le\Delta$. The two tail differences
are $D_1=\Delta-A$ and $D_2=A$, proving \eqref{eq:n2-tail}.
The signed product mass differences at $00,10,01,11$ are
\[
 A-\Delta,\qquad d_1-A,\qquad d_2-A,\qquad A.
\]
Taking half the sum of their absolute values proves \eqref{eq:n2-product}.
The lower bound $\Delta/2$ in both formulas is attained by the
complementary homogeneous pair. The tail equality condition follows
immediately. Product equality is equivalent to $d_1=d_2=A=d$.
Since $d>0$ and
\[
 A=(q_1+d)(q_2+d)-q_1q_2=d(q_1+q_2+d),
\]
this is equivalent to $q_1+q_2=1-d$. Writing $q_1=r$ gives
\eqref{eq:n2-family}, with the displayed interval exactly expressing
feasibility. Conversely every member of that family attains equality;
its homogeneous member is $r=(1-d)/2$.

For strict containment, choose $0<\varepsilon<\min\{d,1-d\}$, set
$d_1=d+\varepsilon$, $d_2=d-\varepsilon$, and take
$p_i=(1+d_i)/2$, $q_i=(1-d_i)/2$. Then $A=d$, so
$T(p,q)=d$ but $\TV(P_p,P_q)=d+\varepsilon>d$.
At $\Delta=0$, ordering forces $p=q$, and every such pair minimizes both
objectives. At $\Delta=2$, the only feasible pair is $p=(1,1)$, $q=(0,0)$.
\end{proof}

\section*{AI assistance}
GPT 6 Sol and Claude Opus 5 assisted with proof discovery, auditing, and
editorial suggestions. The author takes full responsibility for the
correctness of this work.

\begin{samepage}
\begin{thebibliography}{9}
\bibitem{Hoeffding}
W. Hoeffding, \emph{On the distribution of the number of successes in
independent trials}, Ann. Math. Statist. \textbf{27} (1956), no.~3, 713--721.
\href{https://doi.org/10.1214/aoms/1177728178}{doi:10.1214/aoms/1177728178}.

\bibitem{KontorovichTV}
A. Kontorovich, \emph{Minimizing total variation under constraint},
MathOverflow, question 476847, August 13, 2024.
\url{https://mathoverflow.net/q/476847}.

\bibitem{KontorovichMO}
A. Kontorovich, \emph{Poisson binomial conjecture}, MathOverflow,
question 477402, August 22, 2024, with answers by T. Steinke
(August 22, 2024) and W. Sawin (September 24, 2024).
\url{https://mathoverflow.net/q/477402}.

\bibitem{KontorovichTensorization}
A. Kontorovich, \emph{On the tensorization of the variational distance},
Electron. Commun. Probab. \textbf{30} (2025), paper no.~32, 10~pp.
\href{https://doi.org/10.1214/25-ECP680}{doi:10.1214/25-ECP680}.

\bibitem{KontorovichHomogenization}
A. Kontorovich, \emph{TV homogenization inequalities}, preprint (2026),
\href{https://arxiv.org/abs/2601.04079v3}{arXiv:2601.04079v3} [math.PR].

\bibitem{Pitman}
J. Pitman, \emph{Probabilistic bounds on the coefficients of polynomials
with only real zeros}, J. Combin. Theory Ser. A \textbf{77} (1997),
no.~2, 279--303.
\href{https://doi.org/10.1006/jcta.1997.2747}{doi:10.1006/jcta.1997.2747}.
\end{thebibliography}
\end{samepage}
\end{document}
